%%%%%%%%%%%%%%%%%%%%%%%%%%%%%%%%%%%%%%%%%%%%%%%%%%%%%%%
% A template for Wiley article submissions.
% Developed by Overleaf. 
%
% Please note that whilst this template provides a 
% preview of the typeset manuscript for submission, it 
% will not necessarily be the final publication layout.
%
% Usage notes:
% The "blind" option will make anonymous all author, affiliation, correspondence and funding information.
% Use "num-refs" option for numerical citation and references style.
% Use "alpha-refs" option for author-year citation and references style.

\documentclass[num-refs]{wiley-article}
% \documentclass[blind,alpha-refs]{wiley-article}

% Add additional packages here if required
\usepackage{siunitx}

% Update article type if known
\papertype{Original Article}
% Include section in journal if known, otherwise delete
\paperfield{Journal Section}

\title{How do Programmers Maintain Concurrent Code}

% List abbreviations here, if any. Please note that it is preferred that abbreviations be defined at the first instance they appear in the text, rather than creating an abbreviations list.
\abbrevs{ABC, a black cat; DEF, doesn't ever fret; GHI, goes home immediately.}

% Include full author names and degrees, when required by the journal.
% Use the \authfn to add symbols for additional footnotes and present addresses, if any. Usually start with 1 for notes about author contributions; then continuing with 2 etc if any author has a different present address.
%\author[1\authfn{1}]{Feiyue Yu}
\author[1\authfn{1}]{Feiyue Yu}
\author[2\authfn{1}]{Hao Zhong}
\author[1\authfn{2}]{Beijun Shen}
%\author[2]{Author B.~Four}

\contrib[\authfn{1}]{Equally contributing authors.}

% Include full affiliation details for all authors
\affil[1]{School of Software, Shanghai Jiao Tong University, Shanghai, 200240, China}
\affil[2]{Department of Computer Science, Shanghai Jiao Tong University, Shanghai, 200240, China}

\corraddress{Feiyue Yu, School of Software, Shanghai Jiao Tong University, Shanghai, 200240, China}
\corremail{yufeiyue@sjtu.edu.cn}

%\presentadd[\authfn{2}]{Department, Institution, City, State or Province, Postal Code, Country}

\fundinginfo{Funder One, Funder One Department, Grant/Award Number: 123456, 123457 and 123458; Funder Two, Funder Two Department, Grant/Award Number: 123459}

% Include the name of the author that should appear in the running header
\runningauthor{Author One et al.}

\begin{document}

\maketitle

\begin{abstract}
Concurrent programming is pervasive in nowadays software development. Many programmers believe that concurrent programming is difficult, and maintaining concurrency code is error-prone. Although researchers have conducted empirical studies to understand concurrent programming, they still rarely study how programmers maintain concurrent code. To the best of our knowledge, only a recent study explored the modifications on critical sections, and many related questions are still open. In this paper, we conduct an empirical study to explore how programmers maintain concurrent code. We analyze more concurrency-related commits and explore more issues such as the change patterns of maintaining concurrent code than the previous study. We summarize five change patterns according to our analysis on 696 concurrency-related commits. We apply our change patterns to three open source projects, and synthesize three pull requests. Until now, two of them have been accepted. Our results can be useful for programmers to maintain concurrent code and for researchers to implement treating techniques.

% Please include a maximum of seven keywords
\keywords{keyword 1, \emph{keyword 2}, keyword 3, keyword 4, keyword 5, keyword 6, keyword 7}
\end{abstract}

\section{Introduction}
\label{sec:intro}
Many practitioners and researchers believe that the software maintenance phase is one of the most expensive phases, in the life cycle of a software system. Some reports (\emph{e.g.}, \cite{ahn2003software}) claim that the software maintenance phase accounts for almost 80\% of the whole budget. With the maintenance of software, many revision histories are accumulated~\cite{conf/icsm/Borges16}. Based on such revision histories, researchers have conducted various empirical studies to understand how programmers maintain code (\emph{e.g.}, evolution of design patterns~\cite{aversano2007empirical}, fine-grained modifications~\cite{german2006empirical}, and the evolution of APIs~\cite{mcdonnell2013empirical}). These empirical studies deepen our understanding on software maintenance, and provide valuable insights on how to maintain code for programmer.

In recent years, to fully leverage the potential of multi-core CPUs, concurrent programming becomes increasingly popular~\cite{journals/jss/PintoTFFB15}. For example, Pinto \emph{et al.}~\cite{journals/jss/PintoTFFB15} investigated 2,227 projects, and their results show that more than 75\% of their investigated projects employ some concurrency control mechanism. Despite of its popularity, many programmers find that concurrent programming is difficult \cite{journals/corr/McKenney17}, and often introduce relevant bugs in their code~\cite{conf/asplos/LuPSZ08}. As it is difficult to maintain concurrent code, there is a strong need for a thorough empirical study on how programmers maintain such code. Despite of its importance, the topic is still rarely explored. To the best of our knowledge, only a recent study~\cite{conf/sigsoft/GuJSZL15} was conducted to understand how programmers maintain concurrent code. Although the study is insightful and explores many aspects of concurrent programming, it is still incomplete. For example, their study sampled only 25 concurrency-related commits. As another example, they did not analyze the change patterns of maintaining concurrent code. As a result, many relevant questions are still open. For example, are there any patterns, when programmers maintain concurrent code? Indeed, such patterns are useful for programmers when they maintain code. For example, Santos \emph{et al.}~\cite{conf/icsm/SantosAEDV15} have explored the change patterns during software maintenance. Their results show that extracted change patterns can apply in new code locations. However, their study does not touch the change patterns of concurrent code. A more detailed analysis on how programmers maintain concurrent code can have the following benefits:





\noindent
\textbf{Benefit 1.} The results can deepen the knowledge on how to correctly maintain concurrent code. Due to the complexity of concurrent programming, we find that even experienced developers can be confused when they maintain relevant code. For example, Mark Thomas is a member of the Apache Tomcat Project Management Committee\footnote{\url{http://tomcat.apache.org/whoweare.html}}, and senior software engineer at the Covalent division of SpringSource\footnote{\url{https://sourceforge.net/projects/covalent/}}. He contributed more than 10,000 commits to Tomcat. In a commit message, he left the complain as follow:

\lstset{numbers=left, breaklines=true,  basicstyle=\ttfamily\tiny,  xleftmargin=3em, tabsize=2}
%keywordstyle=\color{blue}\bf\ttfamily, language=java,
\begin{lstlisting}
commit a6092d771ec50cf9aa434c75455b842f3ac6c628
Threading / initialisation issues. Not all were valid. Make them volatile anyway so FindBugs doesn't complain.
\end{lstlisting}

\noindent
In this example, we find that even experienced programmers can have problems in understanding their own code changes, when they maintain concurrent code. Our results can reduce such confusions.
%Using existing libraries allows you to  write less code to finish the same work and enjoy the high quality of implementation which is always reliable, strong and fast.


\noindent
\textbf{Benefit 2.} The results can be useful to improve existing tools. For example, Meng \emph{et al.}~\cite{conf/pldi/MengKM11} propose an approach that applies changes systematically based on a given example. With extensions, it can be feasible to apply our extract change patterns to update concurrent code. As another example, our results can be useful to improve existing bug detection approaches \cite{conf/ppopp/SamakR14, conf/sigsoft/EslamimehrP14}.

However, to fulfil the above benefits, we have to overcome the following challenges:

\noindent
\textbf{Challenge 1.} To ensure the reliability of our result, we have to collect many code changes that are related to concurrent programming. It is tedious to manually collect many related code changes for analysis. Tian \emph{et al.}~\cite{tian2012identifying} work on a similar research problem. They propose an approach that identifies bug fixes from commits. Their results show that even advanced techniques can fail to identify many desirable commits.


\noindent
\textbf{Challenge 2.} The changes on concurrent code can be complicated. A recent study~\cite{tufano2016there} show that only 38\% commits are compilable. To analyze code that is not compilable, researchers typically use partial-code tools such as PPA~\cite{DagenaisH08ppa} and ChangeDistiller~\cite{fluri2007change} to analyze commits. However, as partial programs lose information, partial-code tools are imprecise and typically do not support advanced analysis. Furthermore, as we do not know what patterns can be followed, it is difficult to implement a automatic tool. As a result, it is inevitable to take much human effort when we conduct the empirical study.

In this paper, we conducted an empirical study on 98,352 commits that were collected six popular and representative open-source projects. To reduce the effort of manual inspection, we implement a set of tools that collect and identify concurrency-related commits automatically (see Section~\ref{sec:method:tool} for details). With its support, in total, we identified 695 concurrency-related commits, and conducted careful analysis on our identified commits. Based on our results, this paper makes the following contributions:

\begin{itemize}
	\item The first analysis on the change patterns of maintaining concurrent code. Based on our results, we summarize five change patterns, and we present examples to explain each type of our change patterns. We find that following such change patterns, during software maintenance, programmers can modify concurrent code to repair bugs, improve performance, and change functions of their code. Furthermore, we find that maintaining concurrent code is not a one-direction migration. Due to various considerations, programmers can apply seemingly contradicted changes, and even revert their changes. Sometimes, programmers can even make changes, before they fully understand the consequences of their changes.
	\item An application of our change patterns in real code. In particular, we search the latest versions of three projects for chances to apply our change patterns, and synthesize three pull requests according to our change patterns. One of our pull requests is already confirmed and accepted by their programmers. However, our results also reveal that it need much experience and understanding to leverage our change patterns. 
	\item More finding on concurrent program. For example, we reveal the trends of different parallel API classes. As another example, we find that at the granularity of commits, there still exists a correlation between total commits and concurrency-related commits, but it is less significant than at the granularity of code lines. 
\end{itemize}

The rest of paper is organized as follows: Section 2 presents the methodology of our study. Section 3 presents our result and discussion. Section 4 presents related work. Section 5 presents future work and Section 6 concludes.\input


\section{Introduction}
Please lay out your article using the section headings and example objects below, and remember to delete all help text prior to submitting your article to the journal.

\begin{figure}[bt]
\centering
\includegraphics[width=6cm]{example-image-rectangle}
\caption{Although we encourage authors to send us the highest-quality figures possible, for peer-review purposes we are can accept a wide variety of formats, sizes, and resolutions. Legends should be concise but comprehensive – the figure and its legend must be understandable without reference to the text. Include definitions of any symbols used and define/explain all abbreviations and units of measurement.}
\end{figure}

\subsection{Second Level Heading}
If data, scripts or other artefacts used to generate the analyses presented in the article are available via a publicly available data repository, please include a reference to the location of the material within the article.

% Equations should be inserted using standard LaTeX equation and eqnarray environments, not as graphics, and should be set in the main text
This is an equation, numbered
\begin{equation}
\int_0^{+\infty}e^{-x^2}dx=\frac{\sqrt{\pi}}{2}
\end{equation}
And one that is not numbered
\begin{equation*}
e^{i\pi}=-1
\end{equation*}

\subsection{Adding Citations and a References List}

Please use a \verb|.bib| file to store your references. When using Overleaf to prepare your manuscript, you can upload a \verb|.bib| file or import your Mendeley, CiteULike or Zotero library directly as a \verb|.bib| file\footnote{see \url{https://www.overleaf.com/blog/184}}. You can then cite entries from it, like this: \cite{lees2010theoretical}. Just remember to specify a bibliography style, as well as the filename of the \verb|.bib|.

You can find a video tutorial here to learn more about BibTeX: \url{https://www.overleaf.com/help/97-how-to-include-a-bibliography-using-bibtex}.

This template provides two options for the citation and reference list style: 
\begin{description}
\item[Numerical style] Use \verb|\documentclass[...,num-refs]{wiley-article}|
\item[Author-year style] Use \verb|\documentclass[...,alpha-refs]{wiley-article}|
\end{description}

\subsubsection{Third Level Heading}
Supporting information will be included with the published article. For submission any supporting information should be supplied as separate files but referred to in the text.

Appendices will be published after the references. For submission they should be supplied as separate files but referred to in the text.

\paragraph{Fourth Level Heading}
% Here are examples of quotes and epigraphs.
\begin{quote}
The significant problems we have cannot be solved at the same level of thinking with which we created them.\endnote{Albert Einstein said this.}
\end{quote}

\begin{epigraph}{Albert Einstein}
Anyone who has never made a mistake has never tried anything new.
\end{epigraph}

\subparagraph{Fifth level heading}
Measurements should be given in SI or SI-derived units.
Chemical substances should be referred to by the generic name only. Trade names should not be used. Drugs should be referred to by their generic names. If proprietary drugs have been used in the study, refer to these by their generic name, mentioning the proprietary name, and the name and location of the manufacturer, in parentheses.

\begin{table}[bt]
\caption{This is a table. Tables should be self-contained and complement, but not duplicate, information contained in the text. They should be not be provided as images. Legends should be concise but comprehensive – the table, legend and footnotes must be understandable without reference to the text. All abbreviations must be defined in footnotes.}
\begin{threeparttable}
\begin{tabular}{lccrr}
\headrow
\thead{Variables} & \thead{JKL ($\boldsymbol{n=30}$)} & \thead{Control ($\boldsymbol{n=40}$)} & \thead{MN} & \thead{$\boldsymbol t$ (68)}\\
Age at testing & 38 & 58 & 504.48 & 58 ms\\
Age at testing & 38 & 58 & 504.48 & 58 ms\\
Age at testing & 38 & 58 & 504.48 & 58 ms\\
Age at testing & 38 & 58 & 504.48 & 58 ms\\
\hiderowcolors
stop alternating row colors from here onwards\\
Age at testing & 38 & 58 & 504.48 & 58 ms\\
Age at testing & 38 & 58 & 504.48 & 58 ms\\
\hline  % Please only put a hline at the end of the table
\end{tabular}

\begin{tablenotes}
\item JKL, just keep laughing; MN, merry noise.
\end{tablenotes}
\end{threeparttable}
\end{table}

\section*{acknowledgements}
Acknowledgements should include contributions from anyone who does not meet the criteria for authorship (for example, to recognize contributions from people who provided technical help, collation of data, writing assistance, acquisition of funding, or a department chairperson who provided general support), as well as any funding or other support information.

\section*{conflict of interest}
You may be asked to provide a conflict of interest statement during the submission process. Please check the journal's author guidelines for details on what to include in this section. Please ensure you liaise with all co-authors to confirm agreement with the final statement.

\printendnotes

% Submissions are not required to reflect the precise reference formatting of the journal (use of italics, bold etc.), however it is important that all key elements of each reference are included.
\bibliography{sample}

\begin{biography}[example-image-1x1]{A.~One}
Please check with the journal's author guidelines whether author biographies are required. They are usually only included for review-type articles, and typically require photos and brief biographies (up to 75 words) for each author.
\bigskip
\bigskip
\end{biography}

\graphicalabstract{example-image-1x1}{Please check the journal's author guildines for whether a graphical abstract, key points, new findings, or other items are required for display in the Table of Contents.}

\end{document}
